\chapter{资源规划}
\label{lab:A10_resource_planning}

\noindent\textbf{配套文件：}\path{notebook/A10_resource_planning.ipynb}\par
\noindent\textbf{教材关联：}\bookxrefshort{ch:moe}、\bookxrefshort{ch:long-context}、\bookxrefshort{ch:open-models}、\bookxrefshort{ch:resource-budget}、\bookxrefshort{ch:pretraining-scaling}、\bookxrefshort{ch:compute-infrastructure}。

\section*{实验目标}
将显存、计算期限和主机资源分别计量。

\section*{准备及定位}
先阅读本册实验总则，确认Notebook中的依赖与资产单元。原理计算、库接口迁移和真实模型训练可能具有不同资源要求，应分别记录设备、版本及运行范围。

源码定位可从下列函数开始；应结合前置数据与配置单元阅读。
\begin{itemize}
\item \texttt{my\_estimate\_parameter\_breakdown}
\item \texttt{my\_estimate\_inference\_memory}
\item \texttt{my\_estimate\_training\_memory}
\end{itemize}

\section*{操作及观察}
\begin{enumerate}
\item 从参数形状核对参数数目，区分常驻权重与活跃计算。
\item 分别列出推理权重、KV缓存、工作区以及训练梯度和优化器状态。
\item 改变批次、上下文与精度，观察各项预算如何变化。
\item 将GPU数量与期限、CPU、内存、存储及网络需求联合审查。
\end{enumerate}

\section*{结果判读及提交}
提交分项字节表、单位换算和敏感性分析。估计值不是实测峰值；未计碎片与通信不能声称配置一定可运行。

提交时附上所执行单元范围、环境与制品身份、原始结果及失败说明。将未运行项目明确列出，不能用Notebook中已有输出替代本次执行证据。
